% !TEX root = Frozen Braised Napa Cabbage.tex
\documentclass[14pt]{extarticle}
\usepackage[utf8]{inputenc}
% --- REMOVE PAGE NUMBERING ---
\pagestyle{empty}
% --- MARGIN SETUP ---
\usepackage[left=0.75in, right=0.75in, top=0.5in, bottom=0.5in]{geometry}
% --- LIST SETTINGS ---
\usepackage{enumitem}
\setlist[itemize]{leftmargin=1em}
\setlist[itemize]{leftmargin=1em, itemsep=0.2em, label={\raisebox{-0.1ex}{\large\textbullet}}}
\setlist[enumerate]{leftmargin=1.5em, label=\textbf{\arabic*.}, itemsep=0.25em}
\usepackage{parskip}
% --- CONTROL SECTION FONT SIZE ---
\usepackage{titlesec}
\titleformat{\section}{\large\bfseries}{}{0em}{}
\titleformat{\subsection}{\normalsize\bfseries}{}{0em}{}
\titlespacing*{\subsection}{0pt}{0.75em}{0.25em}
\titlespacing*{\section}{0pt}{1em}{0.5em}
% --- FONT SETUP ---
\usepackage{helvet}
\renewcommand{\familydefault}{\sfdefault}
% --- CUSTOM COMMAND FOR TITLE ---
\newcommand{\makeRecipeTitle}[1]{
    {\noindent \Large \textbf{#1}}
    \par \vspace{-0.75em}
}
% --- CUSTOM COMMAND FOR SERVINGS ---
% Usage: \recipeServings{4--5} or \recipeServings{6 (one burrito each)}
% Place directly after the first \makeRecipeTitle on page 1.
\newcommand{\recipeServings}[1]{
    {\noindent \itshape Serves #1 \par}
    \vspace{-0.75em}
}
% --- CUSTOM COMMAND FOR SOURCE CREDIT ---
% Usage: \recipeSource{Creator -- Video Title}{https://...}
% Place just before \end{document}; pins the credit to the bottom of the last page.
\usepackage[hidelinks]{hyperref}
\urlstyle{same}
\newcommand{\recipeSource}[2]{
    \vfill
    {\noindent \footnotesize \textbf{Source:} #1 \\ \url{#2}}
}


\begin{document}

% ================= PAGE 1 =================
\makeRecipeTitle{Frozen Braised Napa Cabbage}

\section*{Ingredients}

\begin{itemize}
      \item Napa Cabbage – 2 medium (any cabbage works here)
      \item Avocado Oil – 2 tbsp (or any high-heat oil)
      \item Ghee – 1 tbsp
      \item Garlic Cloves – 6, peeled
      \item Ginger – 3 knobs, peeled
      \item Jalapeño Pepper – 1 small, deseeded
      \item Chicken Bone Broth – 8 oz (Kettle\&Fire or similar)
      \item Water – 1 cup
\end{itemize}

\section*{Notes}

\begin{itemize}
      \item \textbf{Braise Time:} The longer the braise, the softer the cabbage. A shorter braise leaves more liquid in the pan and a crunchier bite, so cook it to the texture you want.
\end{itemize}

\newpage
\makeRecipeTitle{Frozen Braised Napa Cabbage}
\section*{Instructions}

\begin{enumerate}

      \item \textbf{Prep the Cabbage} \\
            Clean the cabbages and cut each one into quarters through the core, so the leaves stay attached.

      \item \textbf{Sear the Cabbage} \\
            Heat the avocado oil in a large pan over medium-high heat. Sear the cabbage quarters until browned on both cut sides (2--3 minutes per side). Remove from the pan and set aside.

      \item \textbf{Sauté Aromatics} \\
            Add the ghee to the same pan along with the garlic, ginger, and jalapeño. Cook until softened and fragrant.

      \item \textbf{Braise} \\
            Return the cabbage to the pan. Pour in the bone broth and water, cover, and bring to a boil. Uncover and let the liquid reduce for 5 minutes, or longer for softer cabbage.

      \item \textbf{Portion and Freeze} \\
            Transfer to silicone molds and freeze overnight.

\end{enumerate}

\end{document}
